\chapter{Flusso}

\section{Problemi di flusso a costo minimo}

\subsection{Ammissibilit\`a e degenerazione}

Se i flussi associati agli archi in base hanno valore \textbf{non negativo},
si parla di \textbf{soluzione ammissibile}.


Se uno o pi\`u flussi associati agli archi in base hanno valore \textbf{nullo},
si parla di \textbf{soluzione degenere}.


\subsection{Coefficienti di costo ridotto}

Ad ogni arco fuori base \`e assocaito un \textbf{coefficiente di costo ridotto} che misura
la variazione del valore dell'obiettivo al crescere di una unit\`a del flusso associato all'arco fuori base.



\subsection{Condizione di ottimalit\`a}
Se i coefficienti di costo ridotto di tutti gli archi fuori base sono \textbf{non negativi},
in tal caso la soluzione di base attuale \`e \textbf{ottima}.

\subsubsection{Calcolo del costo ridotto}

\begin{enumerate}
  \item Prendere un arco fuori base.
  \item Aggiungere l'arco all'albero di supporto corrispondente alla base attuale.
  \item Considerare l'unico ciclo che si forma con tale aggiunta.
  \item Fissare come verso del ciclo quello dell'arco fuori base.
  \item Calcolare il coefficiente di costo ridotto sommando tra loro tutti i costi relativi agli archi attraversati dal ciclo nel loro stesso verso e sottraendo al risultato i costi degli archi attraversati dal ciclo in senso opposto.
\end{enumerate}


\subsection{Condizione di illimitatezza}
Se l'aggiunta all'albero di supporto/base attuale di un arco fuori base con coefficiente di costo ridotto negativo crea un ciclo orientato, allora il problema di flusso a costo minimo ha obiettivo illimitato.


\subsection{Cambio di base}

Se le condizioni di ottimalit\`a e illimitatezza non sono soddisfatte si procede con il cambio di base
scegliendo un'arco da inserire e un'arco da rimuovere dalla base.

Per gli archi entranti in base bisogna scegliere quelli con coefficiente di costo ridotto negativo.

\subsection{Algoritmo del simplesso su rete}
Una volta completate le operazioni relative a una base e definita una nuova base, l'algoritmo del simplesso su rete ripete le operazioni appena viste sulla nuova base, cioè:



\begin{enumerate}
  \item Verifica di ottimalità.
  \item Verifica di illimitatezza.
  \item Cambio di base.
\end{enumerate}

E itera questo procedimento fino a quando arriva a una base che soddisfi la condizione di ottimalità o quella di illimitatezza.




\section{Problemi di flusso massimo}
Vogliamo determinare la quantità massima di flusso che
partendo dal nodo sorgente si può far giungere fino al nodo
destinazione, tenuto conto dei vincoli di capacità sugli archi
e di quelli di equilibrio nei nodi intermedi.

\subsection{Tagli nella rete}
Considero $U \subset V$ dove $V$ rappresenta l'insieme dei vertici e aggiungo due propriet\`a:
\begin{itemize}
  \item $S \in U$, $S$ \`e il nodo sorgente
  \item $D \not \in U$, $D$ \`e il nodo destinazione
\end{itemize}

Posso separare il grafo di partenza nei due insieme $U$ e $V \setminus U$.

Il taglio \`e l'insieme degli archi che collegano i due insiemi.

\section{Problema di taglio minimo}

\subsection{Archi saturi}
Supponiamo di avere un flusso ammissibile $X = (x_{ij})$ per cui sono soddisfatti i vincoli di equilibrio e quelli di capacità degli archi. Se $x_{ij} = c_{ij}$, diremo che l'arco $(i, j)$ è saturo.

Dove con $c_{ij}$ si intende la capacità dell'arco $(i, j)$.


\subsection{Cammini privi di archi saturi}
Se un cammino dal nodo $S$ al nodo $D$ non presenta nessun arco saturo, il flusso $X$ non è ottimo.

Possiamo aumentare il flusso lungo ciascun arco del cammino di una quantità $\Delta$ definita nel seguente modo:
\[
\Delta = \min_{i=0,...,r} \left[ c_{q_iq_{i+1}} - x_{q_iq_{i+1}} \right] > 0,
\]
senza violare i vincoli di capacità degli archi.

Si noti anche che i vincoli di equilibrio continuano ad essere rispettati, poiché in ogni nodo intermedio del grafo facente parte del cammino, il flusso in entrata aumenta di $\Delta$, ma contemporaneamente aumenta anche il flusso in uscita di $\Delta$.


\subsection{Grafo associato ad un flusso ammissibile}

Definisco un nuovo grafo orientato $G(X) = (V, A(X))$.

Il nuovo grafo ha gli stessi nodi della rete originaria e ha il seguente insieme di archi:

\[
A(X) = \{(i, j) : (i, j) \in A, x_{ij} < c_{ij}\} \cup \{(i, j) : (j, i) \in A, x_{ji} > 0\}
\]

L'insieme $A(X)$ contiene tutti gli archi di $A$ che non sono saturi (l'insieme $A_f(X)$, detto insieme degli archi forward) e tutti gli archi di $A$ lungo cui è stata inviata una quantità positiva di flusso, ma invertiti di verso (l'insieme $A_b(X)$, detto insieme degli archi backward).



\begin{itemize}
  \item Per gli archi \textbf{forward}, il valore $\alpha_{ij}$ rappresenta quanto flusso posso ancora inviare lungo l'arco $(i, j) \in A$ della rete originaria.
  \item Per gli archi \textbf{backward}, il valore $\alpha_{ij}$ rappresenta quanto flusso posso rispedire indietro lungo l'arco $(j, i) \in A$ della rete originaria.
\end{itemize}


\subsection{Aggiornamento del flusso}
Per prima cosa devo trovare il $\Delta$ di un cammino, si deve prendere il valore minore di un cammino.

Dopo di che posso aggionrare il flusso con la seguente formula:
\begin{align}
  x_{ij} = 
  \begin{cases}
    x_{ij} + \Delta  \quad \text{se } (i, j) \in A_f(X) \cap P \\
    x_{ij} - \Delta  \quad \text{se } (j, i) \in A_b(X) \cap P \\
    x_{ij} \quad \text{altrimenti}
  \end{cases}
\end{align}


L'aggiornamento del flusso garantisce che questo sia ancora ammissibile, cioè soddisfa i vincoli di non negatività delle variabili, i vincoli di capacità sugli archi e i vincoli di equilibrio nei nodi intermedi.

Nella rete originaria, incrementiamo di $\Delta$ il flusso lungo gli archi corrispondenti ad archi forward del grafo associato al flusso attuale, appartenenti al cammino $P$. Rispediamo indietro $\Delta$ unità di flusso lungo gli archi che, cambiati di verso, corrispondono ad archi backward del grafo associato, appartenenti al cammino $P$, e quindi decrementiamo il flusso della stessa quantità lungo tali archi. Infine, lungo gli archi che non fanno parte del cammino $P$, il flusso rimane invariato.

Il nuovo flusso è ammissibile e la quantità di flusso uscente da $S$ è anch'essa aumentata proprio della quantità $\Delta$, il che dimostra che il flusso precedente non era ottimo.


\section{Algoritmo di Ford-Fulkerson per il flusso massimo}

\subsection*{Passo 0 - Inizializzazione}
È dato un flusso ammissibile iniziale $X$ (si noti che si può sempre partire con il flusso nullo $X = 0$).

\subsection*{Passo 1}
Si costruisca il grafo $G(X)$ associato al flusso attuale $X$.

\subsection*{Passo 2}
Se non esiste alcun cammino orientato da $S$ a $D$ in $G(X)$, allora STOP: il flusso attuale $X$ è ottimo.
Altrimenti, si individui un tale cammino $P$ e si aggiorni il flusso $X$ come visto in precedenza e si ritorni al Passo 1.

\section{Procedura di etichettatura}
\subsection*{Passo 0}
Si parta con un flusso ammissibile $X$. Si noti che è sempre possibile partire con il flusso nullo $X = 0$.

\subsection*{Passo 1}
Si associ alla sorgente $S$ l'etichetta $(S, \infty)$. L'insieme $R$ dei nodi analizzati è vuoto, ovvero $R = \emptyset$, mentre l'insieme $E$ dei nodi etichettati contiene il solo nodo $S$, ovvero $E = \{S\}$.

\subsection*{Passo 2}
Se $E \setminus R = \emptyset$, il flusso attuale $X$ è ottimo. Altrimenti, si selezioni un nodo $i \in E \setminus R$, cioè etichettato, con etichetta $(k, \Delta)$, ma non ancora analizzato, e lo si analizzi. Per ogni nodo $j \notin E$, cioè non ancora etichettato, e tale che $(i, j) \in A(X)$, si etichetti $j$ con la seguente etichetta:
\begin{align}
  (i, \min(\Delta, c_{ij} - x_{ij}))$ se $(i, j) \in A_f(X)\\
  (i, \min(\Delta, x_{ji}))$ se $(i, j) \in A_b(X)
\end{align}

Quindi si ponga:
\begin{align}
  E &= E \cup \{j \notin E : (i, j) \in A(X)\} \\
  R &= R \cup \{i\}
\end{align}
Se $D \in E$, cioè se la destinazione è stata etichettata, si vada al Passo 3, altrimenti si ripeta il Passo 2.

\subsection*{Passo 3}
Si ricostruisca un cammino orientato da $S$ a $D$ in $G(X)$ procedendo a ritroso da $D$ verso $S$ ed usando le prime componenti delle etichette. A questo punto il cammino trovato
\begin{align}
S \to q_1 \to \ldots \to q_{r-1} \to q_r \to D
\end{align}
con insieme di archi
\begin{align}
P = \{(S, q_1), \ldots, (q_{r-1}, q_r), (q_r, D)\}
\end{align}
è un cammino orientato da $S$ a $D$ in $G(X)$.
Ora possiamo aggiornare $X$ come visto in precedenza, dove il valore $\Delta$ è dato dalla seconda componente dell'etichetta del nodo $D$, e ritornare al Passo 1.



